\section{Data Representation and Preprocessing for Temporal Reinforcement Learning}

The Transformer-based architectures require temporally structured observations rather
than isolated environment states. However, the trading environments provide observations
containing information only for the current timestep:
\[
s_t =
[s_t^{portfolio}, s_t^{market}]
\]
without historical context.

\subsection{Sequence Construction via Rolling Window Buffering}

To enable temporal modeling, a sequence construction mechanism is introduced through the
\texttt{VecSequenceWrapper}. The wrapper maintains a rolling buffer of the most recent
$T$ observations for every parallel environment instance:
\[
S_t
=
[s_{t-T+1}, s_{t-T+2}, \dots, s_t]
\]
where $T=10$ in the experimental configuration.

At each environment step, the oldest observation is removed and the newest state is
appended, producing a sliding temporal window. Consequently, the policy receives a
sequence of past states instead of a single observation.

This mechanism is required only for temporal architectures (Transformers), since MLP and
standard recurrent policies internally manage temporal information differently.

\subsection{Online Rolling Window Normalization}

Since financial variables exhibit substantial differences in scale (prices, holdings,
portfolio values, technical indicators and rewards) and market distributions evolve
continuously over time, training directly on raw observations may lead to unstable
optimization and dominance of high-magnitude features.

To address this issue, a rolling normalization mechanism is introduced through the
\texttt{RollingWindowNorm} wrapper. Unlike static normalization approaches computed
over the entire dataset, the proposed method maintains a fixed-size temporal buffer
containing the most recent observations and computes normalization statistics online.

Importantly, normalization is performed \textbf{independently for each feature
dimension}. Consider an observation vector:
\[
x_t =
[x_t^{(1)}, x_t^{(2)}, \dots, x_t^{(D)}]
\]
where $D$ denotes the feature dimensionality. For each feature $d$, the wrapper keeps
a rolling history over the latest $W$ observations:
\[
\mathcal{B}_t^{(d)}
=
\{x_{t-W+1}^{(d)},\dots,x_t^{(d)}\}
\]
and computes feature-specific statistics:
\[
\mu_t^{(d)}
=
\frac{1}{W}
\sum_{k=t-W+1}^{t}
x_k^{(d)}
\]

\[
\sigma_t^{(d)}
=
\sqrt{
\frac{1}{W}
\sum_{k=t-W+1}^{t}
(x_k^{(d)})^2
-
(\mu_t^{(d)})^2
}
\]

The normalized observation is then obtained element-wise:
\[
\hat{x}_t^{(d)}
=
\frac{
x_t^{(d)}-\mu_t^{(d)}
}
{
\sigma_t^{(d)}+\epsilon
}
\]

Consequently, prices, portfolio quantities and technical indicators are normalized
independently according to their own recent history rather than sharing global
statistics.

Reward normalization follows the same principle using an independent rolling buffer:
\[
\hat{r}_t
=
\frac{r_t-\mu_t^r}
{\sigma_t^r+\epsilon
}
\]
where $\mu_t^r$ and $\sigma_t^r$ are computed over recent rewards.

A key practical advantage of this approach is that it is fully \textbf{online and
self-contained}: the normalization statistics are never stored or fixed after training.
Instead, they are continuously recomputed during both training and evaluation. This
eliminates the need for a separate preprocessing step and avoids the common issue of
mismatched training and test distributions, which is particularly pronounced in
financial markets due to regime shifts and non-stationarity.

This adaptive normalization strategy is therefore especially suitable for trading
environments, where distributional shifts between training and testing data are
expected rather than exceptional. By continuously adapting to recent market conditions,
the wrapper improves stability and robustness of the learned policy under changing
market regimes.